\clearpage\thispagestyle{empty}
\pdfbookmark[0]{Quellen zum Brief über den zweiten Kometen von 1805}{cometletter18060708-sources}
\begin{center}\textbf{Quellen zum Brief über den zweiten Kometen von 1805}\end{center}
\noindent Diese KI-Übertragung in heutiges Deutsch folgt dem Brief Carl Friedrich Gauß’ vom 8. Juli 1806 an Franz Xaver von Zach in \href{https://gdz.sub.uni-goettingen.de/id/PPN236007467}{Gauß’ \textit{Werken}, Band~VI} (Göttingen, 1874), S.~275--277. Der Brief erschien in der \textit{Monatlichen Correspondenz}, Band~XIV, August 1806, Beitrag~XV, \href{https://www.digitale-sammlungen.de/en/view/bsb10538606?page=195}{S.~181}--\href{https://www.digitale-sammlungen.de/en/view/bsb10538606?page=200}{186}. Die Zeitschriftendigitalisate stammen von der Bayerischen Staatsbibliothek.
\par\medskip
Die Werkausgabe enthält die Tagesbruchteile in der Tabelle und liest nach dem Vergleich der Abstände \textit{Aberration}, wo die \href{https://www.digitale-sammlungen.de/en/view/bsb10538606?page=197}{Zeitschrift auf S.~183} \textit{Nutation} liest. Für den Abstand zwischen den äußersten Pallasbeobachtungen nennt die Werkausgabe 71 Tage, die \href{https://www.digitale-sammlungen.de/en/view/bsb10538606?page=200}{Zeitschrift auf S.~186} dagegen 73 Tage.
\par\medskip
Adrien-Marie Legendres \href{https://www.e-rara.ch/zut/content/titleinfo/10838079}{\textit{Nouvelles méthodes pour la détermination des orbites des comètes}} (Paris, Firmin Didot, Jahr XIII--1805) enthalten den Anhang \textit{Sur la Méthode des moindres quarrés}. Übertragung und Satz wurden mit KI erstellt.
\clearpage
